\clearpage\begingroup
\hypertarget{native-331b}{}\HistoricalHead{326}{L}{PRAKTISCHE UND SPHÄRISCHE ASTRONOMIE. BRIEFWECHSEL UND NACHLASS.}
\begin{center}\textsc{Bessel} an \textsc{Gauss}, 10. Juli 1820.\end{center}
\noindent\hspace*{\parindent}\{Ihr lieber Brief vom 28. Juni hat mich auf einen Fehler beim Sirius aufmerksam gemacht. . . . . . . Wenn man diesen Fehler abändert, so geben 23 Beob[achtungen], genau reducirt, \(106^\circ28'33\overset{\prime\prime}{,}57\), und wenn man zwei davon ausschliesst, welche am 7. und 11. März, als ich das Instrument noch zu wenig in Ordnung hatte, gemacht wurden, \(33\overset{\prime\prime}{,}15\). Für \(\alpha\) Virginis finde ich aus 16 Beob[achtungen] \(100^\circ13'5\overset{\prime\prime}{,}50\); für \(\alpha\) Leonis aus 13 B[eobachtungen] \(77^\circ9'24\overset{\prime\prime}{,}01\); für \(\beta\) Leonis aus 14 B[eobachtungen] \(74^\circ25'18\overset{\prime\prime}{,}21\). — Ich bin eben beschäftigt, mehrere Sterne zu reduciren, so dass Sie also noch mehrere Vergleichungspunkte erhalten können und werden.

Inzwischen geht aus Allem hervor, dass Ihre Polardistanzen etwas grösser ausfallen werden, als die meinigen. Es wird sich wohl zeigen, was der Grund davon ist; allein ich bemerke, dass ausser der Biegung noch eine kleine Ursache zur Verschiedenheit vorhanden sein kann, die es vielleicht der Mühe werth wäre ganz zu entfernen, damit wir früher als beim letzten Endresultate \textls[120]{vollkommen} zusammenstimmen. Es ist wohl nicht ganz einerlei, wo man das äussere Thermometer hinhängt. . . . . . .

Ich habe das ununterbrochen trübe Wetter benutzt, einige Beobachtungen zu reduciren, und namentlich habe ich gesucht, die Biegung herauszubringen. Zwei Reihen von Beob[achtungen] geben aber keine gute Uebereinstimmung, indem ich aus der einen diese Biegung im Horizonte \(=2''\) und aus der andern \(=1''\) finde; ich kann noch nicht genau angeben, wie weit dieses ausserhalb der Wahrscheinlichkeit liegt, allein ich setze diese Untersuchung fort und
\clearpage\SourceLeaf{332}\HistoricalHead{327}{R}{BEOBACHTUNGEN AM REICHENBACHSCHEN MERIDIANKREISE.}
\noindent hoffe, wenn nur endlich das Wetter sich ändert, bald der Wahrheit ziemlich nahe zu kommen, wozu die Möglichkeit der Tagbeobachtungen von Wasser förderlich ist. Seit dem 23. Juni habe ich den Polarstern nicht auf diese Art beobachten können, indem nur 3 Culminationen bei heiterem aber unruhigem Wetter vorgefallen sind. Uebrigens glaube ich meine Polhöhe bis auf \(50\overset{\prime\prime}{,}7\) vergrössern und die Declination meiner Tafeln um \(0\overset{\prime\prime}{,}30\) verkleinern zu müssen. Ob diese letzte Aenderung richtig ist, kann ich noch nicht behaupten, da einige Reihen eben so bestimmt eine grössere Verminderung angeben als andere die Richtigkeit der Tafeln, oder gar eine geringe Vermehrung; wobei es mir fast scheint, als ob beide durch den Umstand von einander geschieden werden, dass die Nacht-Culmination in eine Tag-Culmination überging. Es wird überhaupt sehr misslich sein, vor einem Jahre definitive Resultate zu ziehen. . . .
Nach meinen bisherigen Untersuchungen müssen Sie, wenn nicht die Strahlenbrechung einen kleinen Unterschied hervorbringt, bei Ihrem Instrumente die Biegung sehr nahe \(=0\) finden. Ich bin sehr begierig zu hören, was Ihre Untersuchungen wirklich geben werden; doch vielleicht finde ich schon ein vorläufiges Resultat in dem Aufsatze in den G[elehrten] Anz[eigen\begingroup\renewcommand{\thefootnote}{*)}\footnotemark\renewcommand{\thefootnote}{[*)}\footnotetext{Gottingische Gelehrte Anzeigen, Stück 91, 1820 Juni 5; Werke VI, S. 429.]}\endgroup], dessen Sie erwähnen, der aber noch nicht hier ist. . . . . . .\}

\begin{center}\textsc{Bessel} an \textsc{Olbers}, 3. August 1820.\end{center}
\noindent\hspace*{\parindent}\{. . . . . . Sobald ich die Biegung meines Instruments vorläufig \(=1\overset{\prime\prime}{,}5\) bestimmt hatte und die dadurch verbesserten Polardistanzen mit den mir von Göttingen mitgetheilten verglich, erkannte ich, dass das Göttinger Instrument eine weit geringere Biegung haben muss; dass sie aber ganz unmerklich sein sollte, glaube ich kaum. . . . . . .\}

\begin{center}\textsc{Gauss} an \textsc{Bessel}, 17. August 1820.\end{center}
\noindent\hspace*{\parindent}. . . . . . Ich habe seit einiger Zeit die Zenithalsterne beobachtet, die \textsc{Schumacher} in Lauenburg im vorigen Jahre observirt hatte. Jeder Stern ist in beiden Lagen des Instruments 3mal oder darüber beobachtet. 13 Sterne sind jetzt absolvirt und die übrigen 5 oder 6 hoffe ich auch bald zu beendigen. Ich habe jetzt die Reduction auf den Anfang von 1820 gemacht (ganz mit
\clearpage\SourceLeaf{333}\HistoricalHead{328}{L}{PRAKTISCHE UND SPHÄRISCHE ASTRONOMIE. BRIEFWECHSEL.}
\noindent Ihren Elementen). Ich meine \textsc{Schumacher} habe mir gesagt, dass er auch Sie zu Beob[achtungen] dieser Sterne aufgefordert habe. Vielleicht haben Sie also auch bereits die Data zur Vergleichung. Hier meine Resultate:
\begin{center}\begingroup\fontsize{10}{12}\selectfont\setlength{\tabcolsep}{2pt}
\begin{tabular}{@{}lrrrllr@{}}
47 Draconis&\(7^\circ\)&\(38'\)&\(26\overset{\prime\prime}{,}35\)&N.&E[rror] I[ndicis]&\(2'17\overset{\prime\prime}{,}19\)\\
Cephei 2 Hev.&6&50&\(16{,}06\)&&&\(16{,}96\)\\
48 Draconis&6&2&\(55{,}95\)&&&\(17{,}36\)\\
53 Draconis&5&1&\(32{,}34\)&&&\(17{,}06\)\\
33 Cygni&4&29&\(25{,}71\)&&&\(18{,}35\)\\
49 Draconis&3&52&\(25{,}79\)&&&\(17{,}18\)\\
46 Draconis&3&49&\(49{,}51\)&&&\(16{,}62\)\\
51 Draconis&1&35&\(38{,}56\)&&&\(16{,}96\)\\
\(\kappa\) Cygni&1&30&\(37{,}18\)&&&\(17{,}20\)\\
20 Cygni&1&0&\(18{,}21\)&&&\(17{,}54\)\\
\(\iota\) Cygni&0&10&\(46{,}28\)&S.&&\(17{,}06\)\\
\(\iota'\) Cygni praec.&1&25&\(6{,}27\)&&&\(18{,}03\)\\
\(\theta\) Cygni&1&43&\(15{,}97\)&&&\(17{,}79\)\\\cline{2-7}
&\multicolumn{5}{r}{Mittel des C[ollimations-]F[ehlers]}&\(2'17\overset{\prime\prime}{,}33\)
\end{tabular}\par\endgroup\end{center}
\noindent bei gewöhnlicher, d. i. 220maliger Vergr[össerung] des M[eridian]-F[ern]r[ohrs].

\noindent\hspace*{\parindent}Meine gleichzeitigen Circumpolarstern-Beobachtungen gaben den Ort des Pols (so weit sie bisher berechnet sind)
\begin{center}\begingroup\fontsize{10}{12}\selectfont\setlength{\tabcolsep}{2pt}
\begin{tabular}{@{}lllrrrlrrr@{}}
Kreis&in&Ost&\(321^\circ\)&\(29'\)&\(31\overset{\prime\prime}{,}46\)&also P[ol]h[öhe]&\(51^\circ\)&\(31'\)&\(48\overset{\prime\prime}{,}79\)\\
»&»&West&38&25&\(54{,}03\)&&&&\(48{,}64\)\\\cline{7-10}
\multicolumn{7}{r}{Mittel}&\(51^\circ\)&\(31'\)&\(48\overset{\prime\prime}{,}71\)
\end{tabular}\par\endgroup\end{center}
\noindent welches noch die Flexion implicirt. Inzwischen weicht dieselbe nur \(0\overset{\prime\prime}{,}31\) von derjenigen ab, die ich durch die Beob[achtungen] aus dem Wasserhorizont mit den directen abgeleitet, und dies ist eine so kleine Grösse, worüber man bei einem nicht multiplicirenden Instrument mit Verniertheilung wohl \textls[120]{nie} etwas entscheiden kann.
\clearpage\endgroup
